\subsection{Analysis}\label{sec:discussion}

\emph{When the layer acts.} The effect depends on two properties of the
pipeline, not on the dataset: whether the tracker has a continuation step,
and whether the detector's scores match the tracker's thresholds. On MOT17
with well-matched detections, the two-stage trackers BoostTrack, ByteTrack at
floor 0.1, and Hybrid-SORT produced output identical to their hosts, while
the single-stage OC-SORT and PD-SORT gained HOTA on every sequence. When the
stream fits the tracker, the layer stays out of the way, and when the
tracker lacks a step, the layer adds it.

\emph{Self-limiting by design.} In the clean mode the layer keeps the
tracker's thresholds and scores, so a tracker with a low-score step and
well-matched detections receives its normal input. This is the main
difference from the earlier design \aob{}, which lowered the HOTA of the same
strong trackers by 4.15 to 5.89 points (Table~\ref{tab:prior}). The layer
therefore does not need to know in advance whether a tracker needs help; it
reads the answer from the scores.

\emph{A continuation step for single-stage trackers.} Single-stage trackers
drop every detection below one threshold, so a track can be lost when its
object receives low scores for a few frames. The continuation rescue
considers only candidates that overlap a track of the previous frame that no
usable candidate covers, so it targets exactly these gaps. OC-SORT and the
external PD-SORT gained HOTA on 7 of 7 MOT17 sequences through this step.

\emph{Detector replacement.} When a detector is replaced or its score scale
shifts, the fixed thresholds of a tracker no longer match the scores. The
layer adjusts them from the scores without labels. Under halved scores it
restored two trackers from 0 to about 66 HOTA, and with a different detector
it raised the MOTA of two two-stage trackers by 6.0 and 7.9 points. This is
the case the layer was built for: a detector that is swapped, retrained, or
quantized while the tracker stays the same.

\emph{Relation to detector-side control.} AC-MOT \cite{ismail2026acmot} acts
before the detector and chooses its operating point from the scene. \scl{}
acts after the detector and adapts the tracker's input to the scores that the
detector produces. The two layers work on different parts of the pipeline,
so they could be used together; combining them is left for future work.

\emph{Real-time use.} The layer needs no training, no GPU, and no access to
the inside of the detector, and it needs only four contract values from the
tracker's configuration. At run time it reads the candidate list, one small
grayscale image for the motion cue, and the tracks of the previous frame. On
a CPU it takes 2.95 ms per frame, adds 4.5\% to the total time, and makes the
tracker itself faster. \scl{} therefore adds very little delay and is
suitable for real-time use. New trackers are added through their published
contract values and an adapter, without changes to the rules.
